%! Introducing Statistics: What Can We Learn from Data? — Unit 1, Lesson 1.0 — AP Practice
% GRADUAL RELEASE: AP-format practice, assigned but NOT scored — the cover's score column reads
% NA for this component. The lesson's graded practice is the `homework` component that follows
% it at the back of the packet. Shaped like the AP exam: Section I multiple choice in context,
% Section II free response with an interpret-in-context part.
% No spiral-back item: this is the course's first lesson, so there is nothing behind it yet.
% The next-lesson preview is NOT here: it closes the packet, in ../homework/.
% Page lockstep with ap_practice_key rests on \writelines{n} in the blank matching a terse
% \ansline{} in the key — keep key answers short enough not to wrap past the reserved lines.
\documentclass[10pt]{article}
\usepackage{apstats-article}
\usepackage{apstats-boxes}

\begin{document}

\pageheader{Unit 1, Lesson 1.0}{AP Practice}
% NAMESTRIP: no \namedateperiod here — the name row lives on the cover only.

\vspace{0.15in}

% Authored BYTE-IDENTICALLY in ap_practice/ and ap_practice_key/ — this box is what keeps the
% blank and the key the same number of pages.
\boxguard[8]
\begin{remindbox}
\small
This page is \textbf{not required and not scored}. These are AP-format reps: work them for the
practice, and bring whatever stalls you to office hours or study hall. Your \textbf{graded
homework is the last section of this packet}.
\end{remindbox}

\vspace{0.12in}

\boxguard[30]
\begin{scenariobox}[The Adoption Board, Revisited]{navy}
\small
The same county shelter board you worked with in class:

\vspace{8pt}
\renewcommand{\arraystretch}{1.3}
\begin{tabularx}{\linewidth}{@{} l Y Y Y Y @{}}
\rowcolor{sky}
\textbf{Name} & \textbf{Breed} & \textbf{Age (yr)} & \textbf{Weight (lb)} & \textbf{Days on board} \\
\hline
Biscuit & Beagle    & 3 & 24 & 12 \\
Momo    & Terrier   & 1 & 15 & 44 \\
Ranger  & Shepherd  & 5 & 61 &  6 \\
Pepper  & Retriever & 2 & 55 & 21 \\
Sage    & Beagle    & 7 & 27 & 59 \\
Tuck    & Mixed     & 4 & 38 &  9 \\
Willow  & Terrier   & 2 & 17 & 33 \\
Zeke    & Shepherd  & 6 & 70 & 15 \\
\hline
\end{tabularx}
\end{scenariobox}

\vspace{0.14in}

\boxguard[16]
\begin{headlinebox}{goldbg}
\small
\textbf{Section I --- Multiple Choice.} Circle the single best answer. Every answer choice is
about the context, not just the arithmetic --- read them all before choosing.
\end{headlinebox}

\vspace{0.10in}

\begin{enumerate}[leftmargin=1.9em, itemsep=0.14in]

  \item In the adoption board above, which of the following is a \textbf{variable}?
        \begin{enumerate}[label=\textbf{(\Alph*)}, leftmargin=2.2em, itemsep=3pt]
          \item Ranger
          \item Shepherd
          \item 61
          \item Weight (lb)
          \item The eight dogs listed
        \end{enumerate}

  \item Which question about the dogs on this board is a \textbf{statistical question}?
        \begin{enumerate}[label=\textbf{(\Alph*)}, leftmargin=2.2em, itemsep=3pt]
          \item How many dogs are listed on the board?
          \item How old is Willow?
          \item How do the ages of the dogs on this board vary?
          \item Is Tuck a mixed breed?
          \item How many columns does the board have?
        \end{enumerate}

  \item A veterinary clinic records the body temperature, in degrees Celsius, of every dog it
        treats in March. Which statement is \textbf{true}?
        \begin{enumerate}[label=\textbf{(\Alph*)}, leftmargin=2.2em, itemsep=3pt]
          \item The individuals are the temperatures.
          \item The variable is the month of March.
          \item The individuals are the dogs treated in March, and the variable is body temperature.
          \item Body temperature cannot be a variable because it is a measurement.
          \item There is no variable, because only one characteristic was recorded.
        \end{enumerate}

  \item A shelter volunteer says, ``Statistics would be pointless here if every dog were the
        same age, the same weight, and had waited the same number of days.'' The volunteer's
        reasoning depends most directly on the idea that
        \begin{enumerate}[label=\textbf{(\Alph*)}, leftmargin=2.2em, itemsep=3pt]
          \item averages are easier to compute when values differ.
          \item statistics studies characteristics whose values vary across individuals.
          \item a data set must contain at least two variables.
          \item every variable must be measured in units.
          \item counting is not a statistical activity.
        \end{enumerate}

\end{enumerate}

\vspace{0.14in}

\boxguard[16]
\begin{headlinebox}{goldbg}
\small
\textbf{Section II --- Free Response.} Show your reasoning. Every answer must be written
\emph{in the context of the problem} --- that is what earns credit on the AP exam.
\end{headlinebox}

\vspace{0.10in}

\begin{enumerate}[leftmargin=1.9em, itemsep=0.16in, resume]

  \item A school newspaper surveys every student in the senior class. For each senior it records
        the number of hours worked at a paying job last week, and whether the student drives to
        school.

        \vspace{6pt}
        \textbf{(a)} Identify the individuals and both variables in this data set.\\[4pt]
        \writelines{2}

        \vspace{6pt}
        \textbf{(b)} Write one statistical question the newspaper could answer with this data,
        and state precisely what varies in it.\\[4pt]
        \writelines{2}

        \vspace{6pt}
        \textbf{(c)} The paper prints only the sentence \textbf{``Seniors work 11 hours a week.''}
        Identify \textbf{two} pieces of context a reader would need before this number means
        anything.\\[4pt]
        \writelines{2}

        \vspace{6pt}
        \textbf{(d)} A senior reads that sentence and concludes, ``So I should expect to work
        about 11 hours next week.'' Explain why a number that summarizes a whole group does not
        tell you what to expect for one individual in it.\\[4pt]
        \writelines{3}

\end{enumerate}


\end{document}
